\section{Introduction}\label{sec:intro}

\paragraph{The model.}
Consider $N$ electrons of spin one half in the lowest Landau level of a flat torus pierced by
$\Nphi=2N$ flux quanta, so that the filling factor is one half, and let them interact through the
\emph{hollow-core} interaction: the two-body interaction which, for a pair of electrons in the plane, is
the orthogonal projection onto the states of relative angular momentum one. In the language of
pseudopotentials \cite{Hal83} it is the interaction with $V_1=1$ and $V_m=0$ for $m\neq1$. A pair with odd
relative angular momentum has an antisymmetric orbital wave function and therefore a symmetric spin wave
function, so the interaction acts only on spin-triplet pairs; spin-singlet pairs, whose relative angular
momentum is even, cost no energy even at contact, whence the name. The Hamiltonian $H_N$ of this paper
(Definition~\ref{def:hamiltonian}) is the periodization of this interaction on the torus, in second
quantization on the lowest-Landau-level Fock space. It conserves the numbers $\Nup$ and $\Ndn$ of
electrons of each spin, commutes with spin rotations and with the magnetic translations of the torus, and
is nonnegative. Energies are measured in units in which a planar pair of relative angular momentum one has
energy one (Lemma~\ref{lem:units}); statements about gaps that are uniform in $N$ refer to these units.

Haldane and Rezayi \cite{HR88} proposed a spin-singlet incompressible state at filling one half; it has zero energy
for this interaction \cite{RR96,GFN97}. On the torus the zero modes form a degenerate space. At even $N$, in the
sector $\Nup=\Ndn=N/2$, the kernel of $H_N$ is ten-dimensional at every aspect ratio, and at odd $N$ each
of the sectors $\Nup-\Ndn=\pm1$ contains exactly two zero modes. Read and Rezayi \cite{RR96} wrote down the
ten states (the count of ten was first published by Keski-Vakkuri and Wen \cite{KVW93}), and their
completeness is proved in Appendix~\ref{app:kernel}; see Theorem~\ref{thm:kernel}. Whether these zero modes are the ground
states of a gapped phase has been debated. The count of ten was read as the topological ground-state
degeneracy of a non-Abelian phase \cite{KVW93,GFN97,Ino99}; in the description of the Haldane--Rezayi state
by a non-unitary conformal field theory of central charge $c=-2$, the ten states were attributed to a
logarithmic partner of the identity field \cite{GFN97}. Read and Green \cite{RG00} argued instead that the
Haldane--Rezayi state is the critical point of a transition between weak- and strong-pairing phases of
spin-singlet $d$-wave pairing, with gapless bulk excitations, and that both neighbouring gapped phases have
eight states on the torus; general arguments that trial states built from non-unitary conformal field
theories do not describe gapped phases are given in \cite{Rea09a,Rea09b}, and numerical studies of small
systems \cite{CRE19,NPBG23} are consistent with a vanishing gap. To our knowledge the question has not been
settled rigorously.

\paragraph{The result.}
We prove that the zero modes are not separated from the rest of the spectrum by a gap that is uniform in
the number of electrons.

\begin{theorem*}[informal version of Theorem~\ref{thm:main}]
Fix $\lambda>0$, and let $H_N$ be the hollow-core Hamiltonian on the rectangular torus with side lengths
$L_x$, $L_y$, $L_x/L_y=\lambda$, pierced by $\Nphi=2N$ flux quanta. For even $N$ let
$E_1(N)\le E_2(N)\le\cdots$ be the eigenvalues of $H_N$ in the sector $\Nup=\Ndn=N/2$, repeated according to
multiplicity. Then $E_1(N)=\cdots=E_{10}(N)=0$, and $E_{11}(N)\to0$ as $N\to\infty$ through the even
integers.
\end{theorem*}

In other words, there is no $\gamma>0$ such that for infinitely many $N$ the ten-dimensional space of zero
modes is separated by $\gamma$ from the rest of the spectrum in this sector. The theorem follows from a
stronger statement (Theorem~\ref{thm:main}(i)): if for infinitely many even $N$ the $p$ lowest eigenvalues
of the sector lie in an interval of length $\Delta$ and are separated from the rest of the sector by a gap
$\gamma>2\Delta$, with $p$, $\Delta$ and $\gamma$ independent of $N$, then $p$ is divisible by four. We call
such eigenvalues a \emph{uniformly gapped patch} of dimension $p$ at the bottom of the sector: along an infinite
set of $N$, the width of the patch is less than half of its gap to the rest of the sector, uniformly in $N$. The
same condition for a general Hamiltonian and general sectors is Assumption~\ref{ass:gap}, and the term is used
only in this sense; every result stated below for such patches holds also when the width is merely less than the gap
(Remark~\ref{rem:weaker-gap}). Since ten is not divisible by four, the eleventh eigenvalue cannot stay away from zero. At
odd $N$ the same argument gives
$E_3(N)\to0$ in each of the sectors $\Nup-\Ndn=\pm1$, which contain two zero modes each
(Corollary~\ref{cor:oddN}).

\paragraph{The index theorem of Bachmann, Bols, De Roeck and Fraas.}
Our tool is the index theorem of \cite{BBDF21} for gapped quantum lattice systems with a conserved charge.
Suppose that a charge-conserving Hamiltonian of finite range on a lattice has a patch of $p$ eigenvalues at
the bottom of its spectrum, separated from the rest by a gap, and that a locality-preserving,
charge-conserving unitary $U$ leaves the corresponding spectral subspace invariant. Then the charge
transported by $U$ across a cut, averaged over the patch, lies in $p^{-1}\Z$ up to errors that vanish faster
than any power of the size of the system \cite[Theorem~3.2]{BBDF21}. When $U$ is a lattice translation the
transported charge is the average charge of a slab, and the theorem becomes a fractional
Lieb--Schultz--Mattis theorem \cite[Corollary~5.2]{BBDF21}, extending \cite{LSM61,Osh00,Has04,NS07}.

In our setting $U$ is the magnetic translation $\Theta$ by one Landau-gauge orbital, which commutes with
$H_N$, $\Nup$ and $\Ndn$; the charge is the number of up-spin electrons in a band of $N$ consecutive
orbitals; and one application of $\Theta$ moves exactly one orbital's worth of up-spin charge across each of
the two boundaries of the band (Lemma~\ref{lem:theta}). By translation invariance every orbital carries, on
average over a spectral patch of the sector $(N/2,N/2)$, the up-spin charge $(N/2)/\Nphi=1/4$. The index
theorem therefore requires that $p/4$ be an integer up to $\Onf$, hence exactly, since $p$ does not depend
on $N$. For the ten zero modes, $10/4=5/2$ is a half-integer.

\paragraph{Why the index theorem does not apply as written.}
Two features of the lowest Landau level stand in the way.

\emph{The lowest Landau level is not a lattice.} The theorem of \cite{BBDF21} is formulated on a finite graph
with a finite-dimensional space at each vertex, for Hamiltonians that are sums of uniformly bounded terms of
bounded range, and the locality of operators is measured in the graph. The natural basis of the lowest
Landau level of the torus, the Landau-gauge orbitals $\psi_j$, $j\in\Z_{\Nphi}$, is localized in one
direction only: $\psi_j$ is a Gaussian of width one magnetic length about the line $x=2\pi j/L_y$, extended
along it. In this basis $H_N$ is a ring of $\Nphi$ orbitals in which pairs of orbitals at index distance up
to order $L_y$ interact (Theorem~\ref{thm:localform}(c)), and the velocity given by the standard
Lieb--Robinson bound for this chain, measured in orbitals per unit time, grows with $L_y$; a heuristic count,
which we do not make precise, puts it of order $L_y^2$. (In the form of Theorem~\ref{thm:localform}(c), the pair
weights $2d\,e^{-\kappa^2d^2/4}$ of the pairs that contain a given orbital sum to order $\kappa^{-2}$, the pair
operators $\bcen\nu n$ have norms at most of order $\kappa^{-2}$, and the prefactor is $\kappa^3$; so the part of
$H_N$ that does not commute with an even operator at one orbital has norm at most of order $\kappa^{-1}$, and its
range is of order $\kappa^{-1}$ orbitals, which gives a velocity of order $\kappa^{-2}\propto L_y^2$.) The two
boundaries of the band in the argument of
\cite{BBDF21} are $N$ orbitals apart, and at fixed aspect ratio $N$ is of order $L_y^2$, so the time after
which this bound no longer separates the two boundaries does not grow with $N$. The argument of
\cite{BBDF21}, however, filters the dynamics over all times, and it needs the filtered operators to stay
localized near the boundaries (Section~\ref{sec:index}).

\emph{The gap can only be assumed in a sector.} The Fock space of the lowest Landau level contains every
particle number, and at densities below one half the hollow-core Hamiltonian has large kernels; every state
with at most one electron, for instance, is a zero mode. The theorem of \cite{BBDF21} assumes a gap above
the bottom of the spectrum on the whole Hilbert space. For $H_N$ this assumption fails, and a gap can be
assumed at most inside the sector $\Nup=\Ndn=N/2$.

\paragraph{The lowest Landau level as a lattice.}
The main new ingredient of this paper is a way to regard the lowest Landau level of the torus as a
two-dimensional lattice system in which the hollow-core Hamiltonian is a short-range interaction of uniform
size. We construct in Section~\ref{sec:frame} vectors $z_s$ of the lowest Landau level, indexed by the points
$s$ of a rectangular lattice $\Slots$ of \emph{slots} whose spacings are of the order of the magnetic length,
with the following properties.
\begin{enumerate}[(i)]
\item They form a tight (Parseval) frame: $\sum_s\lvert z_s\rangle\langle z_s\rvert$ is the identity on the
lowest Landau level, exactly, for every $N$ (Proposition~\ref{prop:parseval}).
\item Each $z_s$ is localized near the centre $\cen{s}$ of its slot, with Gaussian decay across the
orbitals and decay faster than any power along them, with constants independent of $N$
(Lemma~\ref{lem:localization}).
\item Each orbital $\psi_j$ is a linear combination of the $z_s$ in a single column of slots, two lattice
spacings wide (Lemma~\ref{lem:columns}).
\end{enumerate}
The $z_s$ are windowed Fourier transforms of the orbitals: smooth windows in the orbital index, about one
magnetic length wide, multiplied by Fourier modes along the circumference. The lattice has about twice as
many slots as there are orbitals. A redundant frame rather than an orthonormal basis is natural here: the
lowest Landau level is a band of nonzero Chern number, and for such bands no orthonormal basis made of magnetic
translates of finitely many functions (a system of composite Wannier functions) has finite second moments of the
position \cite{MPPT18}. By (i) the analysis map
$\Sa f=(\ip{z_s}{f})_{s\in\Slots}$ is an isometry, and $c(f)\mapsto\chat(\Sa f)$ extends to an isometric
embedding $\iota$ of the algebra of operators on the lowest-Landau-level Fock space into the algebra of a
lattice fermion system on $\Slots$ with one spin-up and one spin-down mode per slot
(Lemma~\ref{lem:embedding}). The physical Fock space is the subspace of the lattice Fock space in which the
modes orthogonal to the range of $\Sa$ are empty, and $\iota$ intertwines the dynamics and the gauge
transformations generated by $\Nup$ and $\Ndn$.

In this lattice the hollow-core Hamiltonian is local. We show that
$H_N=\int_{\Tor_N}h_R\,\dd R$ exactly, where $h_R=8\pi\sum_\nu B_\nu(R)^*B_\nu(R)$ and $B_\nu(R)$ annihilates
a spin-triplet pair at the point $R$ with relative angular momentum one (Theorem~\ref{thm:localform}); and
that $\iota(h_R)$ is approximated, within $C_k(1+r)^{-k}$ for every $k$, by operators acting only on the
slots within distance $r$ of $R$, with $\norm{h_R}$ and $C_k$ independent of $N$ (Proposition~\ref{prop:slotH}).
Lieb--Robinson bounds \cite{LR72,NS06} hold for lattice fermions with even interactions \cite{NSY18}. We prove
that the dynamics generated by $\iota(H_N)$ carries an operator supported in
a set $X$ of slots to within distance $\rho$ of $X$, up to an error $C_k\norm{A}\abs{X}\rho^{-k}$, whenever
$\rho\ge C_{\rm lc}(1+\abs{t})^\alpha$, for any fixed $\alpha>1$, with constants independent of $N$
(Theorem~\ref{thm:propagation}). The two boundaries of the band are at distance of order $\sqrt N$ magnetic
lengths, so the time up to which this bound keeps them apart now grows with $N$, and the filtered operators
of the argument of \cite{BBDF21} stay localized near the boundaries (Lemma~\ref{lem:Klocal},
Remark~\ref{rem:crossing}). Property (iii) makes the band charges exactly supported in strips of slots, and
conjugation by a band charge preserves strips exactly (Lemma~\ref{lem:bandcharges}); every distance used in
the argument is a distance between strips.

We regard this embedding as a result of independent interest. It turns lowest-Landau-level Hamiltonians
that are integrals of local terms into two-dimensional lattice fermion Hamiltonians to which the locality
theory of lattice systems applies with constants independent of the size of the system. The device itself is not
new: on the sphere, \cite[Section~3]{OAI26} builds a tight frame of the lowest Landau level from windowed
Fourier modes of the orbitals, embeds the physical fermions by its analysis map into the Fock space of a larger set
of fermionic modes, and proves there propagation and time-filter estimates with constants independent of the flux;
it attributes the use of localized frames for the locality of fermions in the continuum to \cite{BDN24}. The frame
constructed here is the analogue on the torus. What the index argument needs from it beyond locality is that each
Landau-gauge orbital is exactly supported in a column of slots, so that the band charges are exactly supported in
strips and conjugation by them preserves strips.

\paragraph{The gap in a sector.}
The assumption of a gap on the whole Hilbert space enters the proof of \cite{BBDF21} in two places: in the
identity $[K,P]=[Q,P]$ for the generator $K$ that dresses the charge $Q$, and in the clustering of the
spectral patch $P$ \cite[Proposition~3.1]{BBDF21}. We show that both remain true when the gap is assumed
only inside a sector, provided they are applied to operators that conserve $\Nup$ and $\Ndn$
(Lemma~\ref{lem:KP} and Proposition~\ref{prop:clustering}), and every operator to which the argument
applies them does. The gap is thus assumed only where it can hold.

\paragraph{The general statement.}
The argument is not specific to the hollow-core interaction. Theorem~\ref{thm:index} is an index theorem for
any sequence of lowest-Landau-level Hamiltonians on rectangular tori which are integrals of terms that
conserve $\Nup$ and $\Ndn$, are uniformly bounded, and are rapidly local in the slot lattice
(Definition~\ref{def:class}); the hollow-core Hamiltonian belongs to this class
(Corollary~\ref{cor:hc-in-class}). For members of the class that are covariant under the magnetic
translation it yields a constraint of Lieb--Schultz--Mattis type (Corollaries~\ref{cor:translation} and
\ref{cor:classwide}): in the sector $\Nup=\Ndn=N/2$, a uniformly gapped patch at the bottom of the spectrum,
in the sense defined above, has dimension divisible by four. In particular no Hamiltonian of the class has in this sector a uniformly
gapped ground-state space of dimension ten. The proof of this constraint, the proof of Theorem~\ref{thm:main}, and Theorem~\ref{thm:kernel} on the zero modes it uses have
been formalized and verified in Lean~4 (Appendix~\ref{app:lean}).

\paragraph{What is excluded, and what is not.}
The theorem excludes a gapped phase whose ground-state space is the space of Haldane--Rezayi zero modes, at
any fixed aspect ratio of a rectangular torus, and uniformly over the aspect ratios of rectangular tori whose
shorter side exceeds a fixed length (Corollary~\ref{cor:uniform}). This is the scenario in which the degeneracy ten of the zero
modes would be the topological ground-state degeneracy of a gapped Haldane--Rezayi phase, as in
\cite{KVW93,GFN97,Ino99}. A uniformly gapped patch at the bottom of the sector, in the sense defined above,
contains the ten zero modes,
so its dimension is at least ten and, being divisible by four, at least twelve; in particular the eight-dimensional ground-state
spaces of the gapped phases adjacent to the Haldane--Rezayi point in the picture of \cite{RG00} are not
ground-state spaces of $H_N$ either. The theorem does \emph{not} exclude a gapped phase whose ground-state
space strictly contains the zero modes and has dimension $12$, $16$, $\dots$: in such a phase
$E_{11}(N)\to0$ would be a splitting inside the ground-state space, not the closing of a gap above it.
The index theorem cannot exclude this, because the conclusions it yields here are divisibility conditions
(Remark~\ref{rem:loophole}). To our
knowledge no such phase has been proposed, and the expectation since \cite{RG00} is that the hollow-core
Hamiltonian has no gap at all; that stronger statement is not proved here. The result is also purely
asymptotic. The constants in our estimates exist but are not explicit, and through the light cone of the
propagation bound of Theorem~\ref{thm:propagation} they are very large; the proof gives neither a rate for $E_{11}(N)\to0$ nor a value of
$N$ beyond which $E_{11}(N)$ is below a prescribed number. At fixed circumference a separate, variational argument with spin
twists bounds $E_{11}(N)$ by $C(L_y)/N$ for every even $N$ and every circumference $L_y$, with $C(L_y)$ explicit
(Theorem~\ref{thm:twist}); this complements the result above and, together with it, gives $E_{11}(N)\to0$ uniformly over
all rectangular tori (Corollary~\ref{cor:alltori}).

\paragraph{Organization.}
Section~\ref{sec:setting} sets up the torus, its lowest Landau level, the Fock space and its sectors, and
the hollow-core Hamiltonian, and quotes the theorem on its zero modes. Section~\ref{sec:frame} constructs
the slot frame and the embedding $\iota$ and defines locality in the slot lattice. Section~\ref{sec:hamiltonian}
proves the local form of the hollow-core Hamiltonian, its uniform size and rapid locality in the slot
lattice and its covariance, introduces the class of Hamiltonians to which the index theorem applies, and
describes the band charges in the slot lattice. Section~\ref{sec:locality} proves the locality estimates: a
Lieb--Robinson bound, propagation from a set with a light cone independent of $N$, the filters, and the
calculus of almost-supported operators. Section~\ref{sec:index} states and proves the index theorem with the
gap assumed in a sector, and Section~\ref{sec:translation} specializes it to the magnetic translation.
Section~\ref{sec:hollowcore} applies the result to the hollow-core Hamiltonian and discusses what it
excludes. Appendix~\ref{app:bbdf} reproduces the operator-theoretic lemmas of \cite{BBDF21} that the proof
uses, Appendix~\ref{app:LR} states the Lieb--Robinson input from \cite{NSY18} and derives the form we use,
Appendix~\ref{app:kernel} proves the theorem on the zero modes, Appendix~\ref{app:twist} proves the variational bound at
fixed circumference, and Appendix~\ref{app:notation} collects the notation.

\paragraph{Conventions.}
Inner products are linear in the second argument. Annihilation operators $c(f)$ are conjugate-linear in
$f$. For $u\in\R$ and $P>0$, $\abs{u}_P$ denotes the distance from $u$ to $P\Z$, and $\pos{u}=\max\{u,0\}$.
Constants denoted $C$, $C'$, $c$ may change from line to line but never depend on $N$; $C_k$ denotes a
constant that may depend on an index $k$. A sequence of operators or numbers $A_N$ is $\Onf$ if for every
$k$ there is $C_k$ with $\norm{A_N}\le C_kN^{-k}$; we write $A\eqN B$ if $A-B=\Onf$, and $x\in\Z+\Onf$ if
$x=x_N$ lies within $\Onf$ of an integer $n_N$ that may depend on $N$.
